% Section 8.8, from lectures/L08/appendix/b_exercises.md.

\section{Exercises}
\label{c8:sec:exercises}

\begin{aside}
\textbf{How to check your work.} Exercise~\ref{c8:ex:binary} and \ref{c8:ex:moore101} are pen and
paper: they are about getting a state machine right \emph{before} it becomes code, and nothing in
them needs VHDL, GHDL or a board.

From Exercise~\ref{c8:ex:mealy101} onwards, every exercise that asks you to write a VHDL module has a
self-checking testbench handed out under \file{lectures/L08/exercises/}, and nothing else;
Exercise~\ref{c8:ex:fastblink} changes a module you already have and is checked by hand. Assembling
the rest of the directory is part of the work: the modules you wrote in earlier chapters
(\code{reset_sync}, \code{button_sync}, \code{sync}, \code{timer}) are copied in from wherever you
wrote them. Every exercise writes out the exact commands it needs.

Exercise~\ref{c8:ex:rx} and \ref{c8:ex:tx} are the book's capstones: the two ends of one and the
same serial link, a receiver and a transmitter. Each composes modules from four different chapters.
Set aside real time for them, and design each state diagram on paper before you write any VHDL. Do
the receiver first: it is the harder of the two, since it has to recover the bit timing it is given
rather than defining it itself.
\end{aside}

% ------------------------------------------------------------------------------------------
\exerciseset{State diagrams and state tables}

\exercise{A machine with binary state ordering}{Circuit}
\label{c8:ex:binary}
Design a new four-state Moore machine, $\{Q_1, Q_2\}$, controlling an LED via a single synchronised,
edge-detected pulse \code{X} from a button (produced exactly like \code{X} in the worked example in
\secref{c8:sec:byhand}).

It differs from the worked example in two ways:
\begin{itemize}
  \item The state ordering is the ordinary binary count
    $00 \to 01 \to 10 \to 11 \to 00$, \textbf{not} the Gray code used in \secref{c8:sec:byhand}.
  \item The LED (\code{Y}) is lit in \textbf{two} states rather than one: whenever $Q_1 = 1$ (that
    is, in both \code{10} and \code{11}).
\end{itemize}

\xpart{a} Draw the state diagram: four states, one transition arrow per state for $X = 1$, and one
self-loop per state for $X = 0$.

\xpart{b} Fill in the complete state table (all eight rows of $\{Q_1, Q_2, X\}$).

\xpart{c} Derive $Q_1^{+}$, $Q_2^{+}$ and \code{Y} with Karnaugh maps, the same way
\secref{c8:sec:byhand} derives them for the Gray-coded version.

\xpart{d} Realise the gate network by hand: two D flip-flops for $\{Q_1, Q_2\}$, gates computing
$Q_1^{+}$, $Q_2^{+}$ and \code{Y}, and build and simulate it in CircuitVerse, with two-flip-flop
synchronisation and edge detection of the button input exactly as in \secref{c8:sec:cv}.

\xpart{e} Verify it in CircuitVerse, against your own state table. There is no testbench for this
machine, and there is not meant to be: the state table you derived in \textbf{b)} \emph{is} the
specification, and checking a circuit against a specification you worked out yourself is the whole
point of the exercise.

Step the clock by hand and check, in this order:
\begin{itemize}
  \item \textbf{The stepping path.} Pulse \code{X} once per step and go a full round,
    $00 \to 01 \to 10 \to 11 \to 00$. Read the expected next state off your table before every edge.
    Read $\{Q_1, Q_2\}$ off the flip-flops' outputs after the edge. All four transitions must match.
  \item \textbf{The holding path.} Step the clock with $X = 0$ in each of the four states and
    confirm that the machine stays put. That is four more checks, and it is the half people skip; a
    machine that steps correctly but does not hold is a machine whose $X'$ terms are wrong.
  \item \textbf{The output.} Check \code{Y} against your table in every state: high in \code{10} and
    \code{11}, low in \code{00} and \code{01}. Do it as a pass of its own, not while you are looking
    at the state bits.
\end{itemize}

Where they differ tells you which gate to look at. A wrong $Q_1^{+}$ shows up only on the
transitions where $Q_1$ should have changed. A wrong \code{Y} shows up with the state sequence still
entirely correct, which is why it is worth a pass of its own.

\note[Hint]Since the state ordering changed from Gray code to binary, expect one transition
($01 \to 10$) where \emph{two} state bits flip at once. Per \secref{c8:sec:byhand} that is harmless
to the machine itself, since nothing samples the state register between edges, and CircuitVerse's
ideal simulation shows you nothing either way. What it does affect is \code{Y}, which here is decoded
combinationally from the state bits: on real hardware \code{Y} can glitch briefly while the two bits
are skewed. Note which transition it is, and what you would do about it if \code{Y} drove something
that cared.

% ------------------------------------------------------------------------------------------
\exerciseset{Sequence detection}

\exercise{A Moore detector for \code{101}}{Reasoning}
\label{c8:ex:moore101}
Design a \textbf{Moore} machine detecting the sequence \code{1, 0, 1} as it arrives one bit per
clock cycle on an input \code{din}, raising its output \code{y} as soon as the hit's third bit has
been seen. Non-overlapping will do here: after a complete hit, start looking for a new
\code{1, 0, 1} from the beginning rather than reusing the closing bits.

\xpart{a} Draw the state diagram: one state per how much of the pattern has matched so far, one
arrow out of every state for \code{din = 0} \emph{and} for \code{din = 1} (a sequence detector is
never idle, it either steps forward or falls back), and mark which state drives \code{y = 1}.

\xpart{b} Fill in the state table, one row per combination of (state, \code{din}).

\xpart{c} Trace it. Pick a six-bit sequence for \code{din} containing at least one hit, and note,
cycle by cycle: the state before the edge, the value of \code{din}, and the output \code{y}.

\xpart{d} The trap in this machine is the arrows back, not the ones forward. For every state, say
what happens on the input that does \emph{not} carry the hit forward: which of them return to the
start, and is there any input that should send you back to a \emph{partially} matched state instead
of all the way to the beginning? Justify your answer for this pattern.

\note[Hint]Count your states before you draw them. ``Nothing matched yet'', ``seen \code{1}'',
``seen \code{10}'', ``seen \code{101}'' is the natural split, and the last of them is what drives
\code{y}.

\textbf{Keep your answer.} Exercise~\ref{c8:ex:mealy101} implements the same detector in VHDL, as a
Mealy machine \textemdash{} which, per \secref{c8:sec:mealy}, needs one state fewer than the one you
just drew. Working out which of your states disappears, and why, is part of that exercise.

% ------------------------------------------------------------------------------------------
\exerciseset{Finite state machines in VHDL}

\exercise{A fourth state in \code{fsm_led}}{VHDL}
\label{c8:ex:fastblink}
Extend \code{fsm_led} with a fourth state, \code{STATE_FAST_BLINK}, blinking at 20~ms instead of
\code{STATE_BLINK}'s 100~ms. Follow the worked example in \secref{c8:sec:vhdl}.

Work on a \textbf{copy} of the module, not on \file{lectures/L08/fsm_led/fsm_led.vhd} itself. First
copy the whole directory somewhere of your own, together with the three modules of yours that it
instantiates:

\begin{shell}
cp -r lectures/L08/fsm_led /tmp/fsm_led_fast
cp lectures/L04/exercises/reset_sync/reset_sync.vhd /tmp/fsm_led_fast     # the three you wrote
cp lectures/L04/exercises/button_sync/button_sync.vhd /tmp/fsm_led_fast   # yourself
cp lectures/L07/exercises/timer/timer.vhd /tmp/fsm_led_fast
cd /tmp/fsm_led_fast
\end{shell}

There is no testbench to run here, but the four-state machine should still analyse, and that is worth
doing after every change: it catches a missing \code{case} branch, or a \code{state_t} value you
added in one process but not the other.

\begin{shell}
ghdl -a --std=93 reset_sync.vhd button_sync.vhd timer.vhd fsm_led.vhd
\end{shell}

\begin{warning}
The reason for working on a copy is worth knowing, since it is a real pitfall and not busywork for
its own sake. \code{fsm_led} has a reference testbench handed out with it,
\file{fsm_led_tb.vhd}, and Exercise~\ref{c8:ex:readfsm} asks you to run it. That testbench checks
the \textbf{three}-state cycle: it presses the button once from \code{STATE_BLINK} and checks that
the LED is steadily on, since the next state in the machine as written is \code{STATE_ON}. In your
four-state version that press lands in \code{STATE_FAST_BLINK}, the LED is still blinking, and the
assertion fails. Nothing in the failure says which of the two exercises it belongs to. Keep the
reference machine intact and that question never arises.
\end{warning}

There is no testbench for the four-state machine. Check it the way \secref{c8:sec:byhand} had you
check a state diagram: predict the state sequence for a series of button presses, then read your
\code{case} branches back and confirm that they agree.

You need to:
\begin{itemize}
  \item Add the new value to \code{state_t}.
  \item Add a branch in every \code{case} statement, in both \code{STATE_PROCESS} and
    \code{LED_PROCESS}.
  \item Handle the two blink rates, either by adding a second timer instance, or by having the
    existing timer's tick count switch on the current state.
\end{itemize}

\exercise{\code{seq_detect_101_mealy}}{VHDL}
\label{c8:ex:mealy101}
Implement the \code{1, 0, 1} detector from Exercise~\ref{c8:ex:moore101} in VHDL, in a new directory
\file{seq_detect_101_mealy} \textemdash{} this time as a \textbf{Mealy} machine.

Start from your state diagram for the Moore machine and convert it:
\begin{itemize}
  \item Your last state, ``seen \code{101}'', exists only to hold the output high for one cycle.
  \item A Mealy machine does not need it: the output condition can read \code{din} directly in the
    state that has already matched \code{1, 0}.
  \item Remove that state, and move its output into the condition for \code{y}.
\end{itemize}

Then follow the \code{type state_t is (...)} and \code{case} pattern from \secref{c8:sec:vhdl} and
\ref{c8:sec:mealy}:
\begin{itemize}
  \item Reuse \code{reset_sync} for \code{reset_n}, exactly as \file{seq_detect_mealy.vhd} does.
    There is no button in this design, so \code{button_sync} is simply not instantiated.
  \item Make sure the assignment of the output \code{y} is a single concurrent statement, outside
    any clocked process, depending on both \code{state} and \code{din}, not just \code{state}.
    Otherwise you have silently built a Moore machine instead, and its output comes one cycle too
    late \textemdash{} which is exactly what the testbench checks.
\end{itemize}

\modulefig[0.56]{lectures/L08/appendix/images/seq_detect_101_mealy.png}

\begin{selfcheck}
\textbf{Self-check.} Name your entity \code{seq_detect_101_mealy}, with the ports \code{clock},
\code{reset_n}, \code{din} (in) and \code{y} (out), declared in that order to match
\code{seq_detect_mealy}. A testbench is in \file{exercises/seq_detect_101_mealy/}; it checks a
non-overlapping \code{101} detector. Your module instantiates \code{reset_sync}, which is not handed
out there since you wrote it in Exercise~\ref{c4:ex:resetsync}: copy your own in.

\begin{shell}
cd lectures/L08/exercises/seq_detect_101_mealy
cp ../../../L04/exercises/reset_sync/reset_sync.vhd .     # the one you wrote in L04
ghdl -a --std=93 reset_sync.vhd seq_detect_101_mealy.vhd seq_detect_101_mealy_tb.vhd
ghdl -e --std=93 seq_detect_101_mealy_tb
ghdl -r --std=93 seq_detect_101_mealy_tb --assert-level=error --stop-time=10ms
\end{shell}
\end{selfcheck}

% ------------------------------------------------------------------------------------------
\exerciseset{Reading the worked machine}

\exercise{Blink rate, testbench and state sequence}{Simulation}
\label{c8:ex:readfsm}
\code{fsm_led} is demonstrated on the DE0-CV (\secref{c8:sec:board}). This exercise checks that you
can account for what you saw, and it comes with a reference testbench you can run yourself. It
reuses your two synchronisers from \chapref{c4:ch} and your \code{timer} from \chapref{c7:ch}, and
none of them is handed out, so copy them in first:

\begin{shell}
cd lectures/L08/fsm_led
cp ../../L04/exercises/reset_sync/reset_sync.vhd .     # the three you wrote yourself
cp ../../L04/exercises/button_sync/button_sync.vhd .
cp ../../L07/exercises/timer/timer.vhd .
ghdl -a --std=93 reset_sync.vhd button_sync.vhd timer.vhd \
                 fsm_led.vhd fsm_led_tb.vhd
ghdl -e --std=93 fsm_led_tb
ghdl -r --std=93 fsm_led_tb --assert-level=error --stop-time=10ms
\end{shell}

\xpart{a} Work out the blink rate from the source, not from the demonstration. \code{fsm_led}
instantiates \code{timer} with \code{TIMER_TICK_COUNT} defaulting to \code{5_000_000}, on a
\code{50 MHz} clock. How long is one timeout period? \code{STATE_BLINK} toggles the LED at every
timeout. What \textbf{blink} frequency does that give, and why is it half of what you might first
write down?

\xpart{b} The testbench overrides \code{TIMER_TICK_COUNT} with a much smaller value than the
default. Find that override in \file{fsm_led_tb.vhd}. Explain why a testbench cannot reasonably use
the real value, and roughly how long a simulation of one real blink period would have to run.

\xpart{c} Account for the state machine's behaviour:
\begin{itemize}
  \item All three states are reachable in both directions. Trace the presses of \code{button_n(0)}
    and \code{button_n(1)} that visit every state and return to \code{STATE_OFF}.
  \item Press both buttons during the same clock cycle. What does \code{fsm_led} do, and which two
    lines in the architecture decide that?
  \item \code{timer_enable} is driven from \code{state}, so the timer runs only in
    \code{STATE_BLINK}. Note that \code{timer} is \textbf{frozen} rather than cleared when it is
    disabled (\secref{c7:sec:module}). What does that mean for the moment you return to
    \code{STATE_BLINK} after leaving it, and would it be better or worse to hold
    \code{timer_enable} permanently high?
\end{itemize}

% ------------------------------------------------------------------------------------------
\exerciseset{Capstone: a serial link}

\exercise{\code{uart_rx8}: the receiver}{VHDL}
\label{c8:ex:rx}
This is the first half of the book's capstone, and it composes four chapters into one design.
Exercise~\ref{c8:ex:tx} builds the other end of the same serial link, so do not stop here. Nothing in
it is a new idea; it is all about wiring together modules you already have, around a state machine
you design yourself:
\begin{itemize}
  \item \chapref{c4:ch}'s \code{reset_sync}, synchronising reset.
  \item \chapref{c5:ch}'s \code{sync}, synchronising the \code{rx} line, which comes from somebody
    else's crystal and is asynchronous in exactly the sense \chapref{c4:ch} means.
  \item \chapref{c7:ch}'s \code{timer}, generating the oversampling tick.
  \item This chapter's state machine, framing all of it and doing its own shifting and bit counting.
    You built that as a module in \chapref{c6:ch}; here it is four lines inside a machine that
    already has the counters it needs, which is the more common way a receiver is written.
\end{itemize}

Write an entity named \code{uart_rx8}: a receiver for one byte of asynchronous serial data, in the
format a UART sends.

\task{The line's format}
The \code{rx} line carries frames, and between frames it idles:
\begin{itemize}
  \item \textbf{Idle}: the line is held high.
  \item \textbf{Start bit}: the line goes low for exactly one bit period. That is what marks the
    beginning of a frame; there is no separate clock line, which is what ``asynchronous serial''
    means.
  \item \textbf{Eight data bits}, most significant first, one bit period each.
  \item \textbf{Stop bit}: the line goes back high for one bit period.
\end{itemize}

\lecfig[0.78]{lectures/L08/appendix/images/uart_frame.png}{A UART frame: a line idling high, a start
bit, eight data bits with the most significant first, and a stop bit, with the receiver sampling
mid-bit.}{c8:fig:uartframe}

The frame above carries \code{10110010}, which is the first byte the testbench sends. The arrows are
the eight moments where the receiver picks up a \textbf{data bit}: one and a half bit periods after
the falling edge, and then once per bit period thereafter. It looks at the line in between as well
\textemdash{} after the falling edge, mid start bit to check that it is still low, and mid stop bit
\textemdash{} but only the eight arrows contribute bits to the received byte.

The receiver has to recover the bit timing from the start bit's falling edge alone. That is the
whole problem this exercise solves, and it is why the timer is here.

\begin{aside}
\textbf{One deliberate departure from a real UART.} A real UART sends its data bits
\textbf{least} significant first, as \secref{c6:sec:sipo} points out. This exercise sends them most
significant first:
\begin{itemize}
  \item the framing and the timing, which are what these two exercises are actually about, are
    identical either way.
  \item it is the direction \secref{c6:sec:sipo}'s SIPO idiom shifts in, so the receiver's shift is
    one line, and the transmitter in Exercise~\ref{c8:ex:tx} is its mirror image.
  \item the testbenches handed out send and expect MSB first to match, so the pair is consistent
    with itself.
  \item to talk to a real UART you would reverse the shift direction to \secref{c6:sec:shift}'s
    mirrored form and change nothing else. Worth knowing before you connect this to a computer's
    serial port and wonder why every byte arrives bit-reversed.
\end{itemize}
\end{aside}

\task{Oversampling: sixteen ticks per bit}
The receiver has to find the middle of a bit whose beginning it can only infer. The technique is
called \textbf{oversampling}: run \code{timer} at \textbf{sixteen times} the baud rate, so that there
are sixteen ticks per bit period, and count them.
\begin{itemize}
  \item The start bit's falling edge clears the tick count to \code{0}. The timer runs continuously,
    before, during and after a frame; it is the \emph{count} that is cleared, never the timer.
  \item \textbf{Tick 8} is half a bit later, in the middle of the start bit. Check that the line is
    \emph{still low} there. If it has gone high it was a glitch and no start bit, and the receiver
    returns to idle. That recheck is what keeps noise on an idle line from being read as data.
  \item From there, every \textbf{sixteenth} tick is the middle of the next bit: sample the line and
    shift it in. Eight of them are the data bits, and the ninth is the stop bit.
\end{itemize}

Sampling mid-bit instead of at the edges is what gives the receiver its timing margin. The
transmitter's clock is a different physical crystal from yours, so the two drift apart over a frame;
a receiver sampling near the boundaries would read the wrong bit as soon as they did. Sixteen ticks
is the conventional choice, and the transmitter in Exercise~\ref{c8:ex:tx} runs on the same tick.

\task{The entity}

\begin{booktable}{@{}p{11em}p{4.4em}p{3.4em}L@{}}
  \thead{Generic} & \thead{Type} & \thead{Default} & \thead{Description} \\
  \midrule
  \code{OVERSAMPLE_TICK_COUNT} & \code{natural} & \code{325} & The tick count for \code{timer} for
    \textbf{one sixteenth} of a bit period. 9600 baud at \code{50 MHz}; part \textbf{b)} asks you to
    derive it. \\
\end{booktable}

\begin{booktable}{@{}p{5.6em}p{3.4em}p{10.8em}L@{}}
  \thead{Port} & \thead{Dir.} & \thead{Type} & \thead{Description} \\
  \midrule
  \code{clock} & in & \code{std_logic} & System clock. \\
  \code{reset_n} & in & \code{std_logic} & Asynchronous, active low; goes through
    \code{reset_sync}, as always. \\
  \code{rx} & in & \code{std_logic} & The serial line, straight from a pin and therefore
    \textbf{asynchronous}. Synchronise it to \code{rx_s2} before anything else reads it. \\
  \code{data_out} & out & \code{std_logic_vector(7 downto 0)} & The received byte, valid during the
    cycle where \code{byte_valid} is high. \\
  \code{byte_valid} & out & \code{std_logic} & A one-cycle pulse: a complete, correctly framed byte
    is on \code{data_out}. \\
  \code{frame_err} & out & \code{std_logic} & A one-cycle pulse: the stop bit was not high, so the
    frame is broken and no byte is handed over. \\
\end{booktable}

\xpart{a} Design the state machine \textbf{on paper first}, exactly as \secref{c8:sec:byhand} had you
do. Four states:
\begin{itemize}
  \item \code{STATE_IDLE}, waiting for the line to go low.
  \item \code{STATE_START}, checking at tick 8 that the line is still low.
  \item \code{STATE_DATA}, sampling one bit every sixteenth tick, eight times.
  \item \code{STATE_STOP}, sampling the stop bit sixteen ticks after the last data bit and choosing
    between \code{byte_valid} and \code{frame_err}.
\end{itemize}

Draw it before you write any VHDL, and on the diagram:
\begin{itemize}
  \item Give every state its \textbf{self-loop}: the arrow for ``the tick I am waiting for has not
    come yet, stay put''. That is the arrow that usually goes missing, and a state without it falls
    through to wherever \code{else} happens to point.
  \item Write beside every state \textbf{how many ticks it waits for}, so that the tick bookkeeping
    is on the paper instead of in your head.
  \item Mark where \code{ticks} is \textbf{cleared}. The timer itself never stops, so those clears
    are the only places where the receiver decides what ``the beginning of a bit'' means.
  \item Label the arrow out of \code{STATE_IDLE} with an \textbf{edge}, not a level: \code{rx_s2}
    low \emph{and} high at the previous tick. That takes one bit of history, so note down where you
    intend to keep it.
\end{itemize}

\xpart{b} Work out the timing before you implement it:
\begin{itemize}
  \item The DE0-CV's clock is \code{50 MHz} and 9600 baud means 9600 bits per second. How many clock
    cycles is one bit period? How many is one sixteenth of one?
  \item \code{timer} pulses every \code{TICK_COUNT + 1} clock cycles (\secref{c7:sec:concept}).
    Which \code{OVERSAMPLE_TICK_COUNT} does that give you? The exact answer is not an integer, so
    there are two candidates either side of it: work out the error each of them leaves, and say
    which of them the default above is and why.
  \item That error is a fraction of a percent, and it accumulates over a frame. How far has your
    sampling point drifted from the middle of the stop bit ten bit periods after the start edge,
    measured in oversampling ticks? How many ticks of margin did it have to begin with?
  \item And now the reason the answer is comfortable: you resynchronise at every start bit. What
    would have to be true of the frame length for that to stop being enough?
\end{itemize}

\xpart{c} Implement \code{uart_rx8}. It instantiates three modules and does the rest itself:
\begin{itemize}
  \item \code{reset_sync} for reset. Everything else in the design takes its \code{reset_s2_n}.
  \item \textbf{\code{sync}}, the module you wrote in Exercise~\ref{c5:ex:sync}, on the \code{rx}
    pin, giving the \code{rx_s2} the state machine reads. Instantiate it with \code{SIZE = 1} and
    \code{PRESET = '1'}: an idle line is high, and a synchroniser coming out of reset reading
    \code{'0'} would look to this machine like a start bit. \code{rx} comes from a pin driven by a
    transmitter with its own crystal, so it is asynchronous in exactly the sense \chapref{c4:ch}
    means, and every rule from there applies unchanged. Note that it takes a one-element
    \code{std_logic_vector} on each side, so you declare a pair of those and pick \code{rx_s2} out
    of the output.
  \item \code{timer}, with \code{OVERSAMPLE_TICK_COUNT} passed straight through, \code{enable} tied
    to \code{'1'}, and its \code{timeout} as your only notion of elapsed time. Call it
    \code{baud_tick}. The timer free-runs: it is never stopped, never cleared, and no state ever
    touches it. What the machine rephases instead is its own counter \code{ticks}.
  \item No shift register module. The machine needs a bit counter anyway, so it shifts the line
    straight into a signal \code{frame} itself, one line of VHDL per sampled bit.
\end{itemize}

\task{The state machine, state by state}
Two counters: \code{ticks} counting \code{0} to \code{15} within a bit, and \code{bit_count} counting
the eight data bits. Everything below happens \textbf{only when \code{baud_tick = '1'}}; at every
other clock edge the machine stands still. Since the timer free-runs, a tick is always on its way, in
every state including idle.

\begin{booktable}{@{}p{6.6em}p{7em}p{10em}L@{}}
  \thead{State} & \thead{Waits for} & \thead{Then} & \thead{To} \\
  \midrule
  \code{STATE_IDLE} & falling edge on \code{rx_s2} & clear \code{ticks} & \code{STATE_START} \\
  \code{STATE_START} & \code{ticks = 8} & \code{rx_s2} still low? clear \code{ticks} &
    \code{STATE_DATA} \\
   & & \code{rx_s2} high? it was a glitch & \code{STATE_IDLE} \\
  \code{STATE_DATA} & \code{ticks = 15} & shift \code{rx_s2} into \code{frame}, clear \code{ticks} &
    \code{STATE_DATA} until 8 bits, then \code{STATE_STOP} \\
  \code{STATE_STOP} & \code{ticks = 15} & \code{rx_s2} high? \code{data_out <= frame}, pulse
    \code{byte_valid} & \code{STATE_IDLE} \\
   & & \code{rx_s2} low? pulse \code{frame_err}, discard the byte & \code{STATE_IDLE} \\
\end{booktable}

There is no ``done'' state. The stop bit's sample is where the frame is judged, so that is where the
byte is handed over, and the machine goes straight back to idle.

\task{Why the idle state waits for an edge instead of a level}
Keep one tick of history for \code{rx_s2}, and leave the idle state only at the transition from high
to low. Testing the level instead, ``if \code{rx_s2} is low, start a frame'', looks equivalent and
handles everything except the one case that matters.

Consider a \textbf{break}: the line held low for far longer than a frame, which is what a transmitter
losing power looks like. A level test rearms the moment the previous frame's stop bit is judged, so
the receiver spends the whole break marching through frame after frame of all zeros. Each ends on a
low stop bit and is correctly rejected, and that part is fine. The problem is the frame in the air
when the break \emph{ends}: its stop bit is sampled after the line has gone back high, so it is a
perfectly well-formed frame as far as the receiver can tell, and a byte made of nothing is handed
over as real data. The edge test does not rearm at all during the break, since after the first
falling edge there is no second one until the line has recovered.

That is the whole difference, it costs one flip-flop, and it is why a receiver that ``works'' on
clean data can still hand over rubbish the first time a cable is pulled out.

\task{The shift}
Eight data bits arrive most significant first, so the SIPO idiom from \secref{c6:sec:sipo} handles it
with no index arithmetic at all:

\begin{vhdlcode}
frame <= frame(6 downto 0) & rx_s2;
\end{vhdlcode}

\code{bit_count} then only counts \emph{how many} bits have arrived, never \emph{where} they land.
Storing with \code{frame(bit_count) <= rx_s2} instead puts the first-arrived bit in bit 0, which
reverses the byte: that is the least-significant-first order a real UART uses, and it would be the
right choice if this receiver had to talk to one. Here the MSB comes first on the line, so the shift
is the one that matches.

Four things decide whether this works:
\begin{itemize}
  \item \textbf{Everything is conditioned on \code{baud_tick}.} A state acting on an ordinary clock
    edge runs sixteen times per tick and about five thousand times per bit.
  \item \textbf{Every state needs its ``not yet'' path.} If the tick you are waiting for has not
    arrived, or the count has not reached its target, the machine must \textbf{stay where it is}.
    Sending it somewhere else in the \code{else} is by far the most common way this design fails, and
    it fails completely rather than subtly: the machine ends up bouncing back and forth between two
    states and never reaches the third.
  \item \textbf{Watch the off-by-one in the counter.} With \code{ticks} starting at \code{0} and
    incrementing at every tick, the $n$th tick is \code{ticks = n - 1} at the moment you test it.
    Whether you compare before or after the increment decides which tick you actually sample on, and
    being one or two ticks late is survivable while eight is not.
  \item \textbf{The timer free-runs, so your phase comes from \code{ticks}, never from the timer.}
    The start edge arrives when the transmitter sends it, which is somewhere inside a tick period the
    receiver did not choose, so the first tick after that edge can come anywhere from a whole tick to
    almost no time later. Clearing \code{ticks} on the way out of the idle state is what rephases the
    sampling against the start bit. The leftover fraction of a tick is the price of never stopping
    the timer, and it is cheap: one sixteenth of a bit, against the half bit of margin part
    \textbf{b)} asks you to work out.
\end{itemize}

\xpart{d} Run the testbench. It needs four design files, the subcomponents first, and none of the
three subcomponents is handed out in the exercise directory: all three are modules you wrote.

\begin{shell}
cd lectures/L08/exercises/uart_rx8
cp ../../../L04/exercises/reset_sync/reset_sync.vhd .     # the three you wrote yourself
cp ../../../L05/exercises/sync/sync.vhd .
cp ../../../L07/exercises/timer/timer.vhd .
ghdl -a --std=93 sync.vhd reset_sync.vhd timer.vhd uart_rx8.vhd uart_rx8_tb.vhd
ghdl -e --std=93 uart_rx8_tb
ghdl -r --std=93 uart_rx8_tb --assert-level=error --stop-time=10ms
\end{shell}

It plays transmitter, and checks:
\begin{itemize}
  \item two clean frames, arriving as the right byte, with no \code{frame_err}.
  \item an idle line, before, between and after them, producing nothing.
  \item a frame whose stop bit is \textbf{low}: \code{frame_err} must pulse and no byte may be handed
    over.
  \item the line held low for far longer than a frame, which is a \textbf{break}: every frame the
    receiver thinks it sees there ends on a low stop bit, so none of them is a byte.
  \item a frame where every data bit carries its value only in the \textbf{middle} of the bit period
    and the opposite level either side. A receiver sampling near tick 8 reads it perfectly; one
    sampling near a boundary reads the neighbouring bit.
  \item that \code{byte_valid} and \code{frame_err} are each never high for two cycles in a row.
\end{itemize}

\note[Hint]Build it in stages and run the testbench after each, rather than writing all four states
and debugging them at once. First get the machine to leave the idle state on a start bit and reach
\code{STATE_DATA}; if it never gets there it is the ``not yet'' path to look at. Then sample one data
bit and check that it is the right bit before you worry about the other seven. Then the stop bit and
the two pulses.

\xpart{e} Now break it deliberately, and think about what the result proves and does not prove:
\begin{itemize}
  \item Remove the stop bit state, and hand the byte over as soon as the eighth data bit is in.
    Predict what should go wrong, then run it. Two checks catch this, and it is worth knowing which
    fires first: the frame with the low stop bit no longer raises \code{frame_err}, and it hands over
    a byte that should have been discarded.
  \item Now the other half of the lesson. Passing is the weaker signal, so find a sabotage the
    testbench does \textbf{not} catch. Here is one: \textbf{remove the mid-start-bit recheck}, so
    that \code{STATE_START} moves on to \code{STATE_DATA} at tick 8 without confirming that the line
    is still low. The testbench never puts a glitch on an idle line, so nothing notices, and on real
    hardware a single spike of noise starts a phantom frame. Read \file{uart_rx8_tb.vhd} and say what
    you would add to catch it.
  \item Here is a second: \textbf{sample at tick 3 instead of tick 8.} Work out why the testbench
    still passes, and what you have given away. The answer is in the margin calculation in part
    \textbf{b)}.
  \item And a third, which is the most important of the three: \textbf{remove the synchroniser for
    \code{rx}} and read the pin directly. Every check still passes, and it always will, whatever is
    added to the testbench. Say why, in terms of what a simulator models and what it does not, and
    then reread \secref{c4:sec:mtbf}. This is the one kind of bug in the book that testing cannot
    reach, which is exactly why \chapref{c4:ch} gives you a rule to follow rather than a symptom to
    look for.
  \item In a few sentences: what does ``passes its testbench'' actually promise, and what does it
    not?
\end{itemize}

\modulefig[0.74]{lectures/L08/appendix/images/uart_rx8.png}

\begin{selfcheck}
\textbf{Self-check.} Name your entity \code{uart_rx8}, with the generic
\code{OVERSAMPLE_TICK_COUNT} (\code{natural}), the inputs \code{clock}, \code{reset_n}, \code{rx},
and the outputs \code{data_out} (\code{std_logic_vector(7 downto 0)}), \code{byte_valid} and
\code{frame_err}, declared in that order. Its testbench is in \file{exercises/uart_rx8/}, and nothing
else: \textbf{copy in your own \code{reset_sync}, \code{sync} and \code{timer}}, as part \textbf{d)}
shows. A project directory contains everything it is built from, and assembling it is part of the
exercise.
\end{selfcheck}

\exercise{\code{uart_tx8}: the transmitter}{VHDL}
\label{c8:ex:tx}
The other half of the link: write \code{uart_tx8}, a transmitter sending a byte in the same format
Exercise~\ref{c8:ex:rx} receives.

A transmitter is the easier of the two, and the reason is worth understanding before you start:
\begin{itemize}
  \item The receiver has to \emph{recover} the bit timing from a start edge it did not control,
    which is why it sampled mid-bit and why its timer ran at one sixteenth of a bit period.
  \item The transmitter \textbf{defines} the timing. Nothing has to be recovered, so the timer runs
    at a whole bit period and the machine simply holds every bit on the line for one timeout.
\end{itemize}

It sends the most significant bit first, which matches the receiver in Exercise~\ref{c8:ex:rx}
rather than a real UART, for the reason given there.

\task{The entity}

\begin{booktable}{@{}p{9em}p{4.4em}p{3.4em}L@{}}
  \thead{Generic} & \thead{Type} & \thead{Default} & \thead{Description} \\
  \midrule
  \code{BIT_TICK_COUNT} & \code{natural} & \code{5207} & The tick count for \code{timer} for a
    \textbf{whole} bit period, 9600 baud at \code{50 MHz}. Sixteen of Exercise~\ref{c8:ex:rx}'s
    oversampling ticks, give or take the rounding its part \textbf{b)} asks you to work out, so that
    a \code{uart_tx8} and a \code{uart_rx8} built with the matching defaults talk at the same
    rate. \\
\end{booktable}

\begin{booktable}{@{}p{5.4em}p{3.4em}p{10.8em}L@{}}
  \thead{Port} & \thead{Dir.} & \thead{Type} & \thead{Description} \\
  \midrule
  \code{clock} & in & \code{std_logic} & System clock. \\
  \code{reset_n} & in & \code{std_logic} & Asynchronous, active low; through \code{reset_sync}, as
    always. \\
  \code{data_in} & in & \code{std_logic_vector(7 downto 0)} & The byte to send. \\
  \code{send} & in & \code{std_logic} & A request to send \code{data_in}. Ignored while the
    transmitter is busy. \\
  \code{tx} & out & \code{std_logic} & The serial line. \textbf{Idles high}, including during
    reset. \\
  \code{busy} & out & \code{std_logic} & High from the moment a frame starts until the stop bit is
    over. \\
\end{booktable}

\xpart{a} Draw the state machine first. Four states are the natural split: idle, start bit, data
bits, stop bit, and each holds the line at a defined level for one bit period.
\begin{itemize}
  \item Mark what drives \code{tx} in every state. Note that \code{tx} never depends on \code{send}
    or any other primary input: it is the state, plus the shift register's current output. That makes
    this a Moore machine, like everything else in this book.
  \item Mark where \code{busy} comes from. It is a single comparison against the state.
  \item Decide where the byte is captured. It has to be latched when the frame starts, not read
    continuously out of \code{data_in}: the producer is free to change \code{data_in} the moment
    \code{busy} goes high, and a real one will.
\end{itemize}

\xpart{b} Shift the byte out most significant bit first, in the machine itself, the mirror image of
what the receiver does: latch \code{data_in} into a signal \code{frame} when the frame starts, drive
\code{tx} from the top bit of \code{frame} while in the data state, and shift \code{frame} one place
up per bit period, counting the eight bits going out.

You wrote this as a module in Exercise~\ref{c6:ex:shiftregs}, and \code{piso8} is worth rereading for
the idiom's sake. Instantiating it here would work, but the machine already owns the bit counter a
PISO register needs, so the shift becomes a single line inside the state where it belongs. That is
how both halves of this link are built, and the same trade-off comes back in \chapref{c14:ch}, where
\code{tx_shift_reg} is a module of its own precisely because \code{can_controller} does \emph{not}
already own the counter it needs.

\xpart{c} Implement it, with \code{reset_sync} instantiated for reset and \code{timer} with
\code{BIT_TICK_COUNT} passed straight through.

Here the timer's \code{enable} \textbf{is} driven from the state, high only while a frame is in the
air, and that is the one place where this design deliberately differs from Exercise~\ref{c8:ex:rx}.
The receiver has to let its timer free-run because it does not control when a frame starts: it can
only rephase its own tick count once the start edge has arrived. The transmitter has the opposite
problem. It \emph{decides} when the frame starts, so it can start the timer at that moment and get a
first bit period exactly as long as all the others. Let it free-run here instead and the start bit is
short by precisely as much of a tick period as happened to remain when \code{send} arrived, which is
a defect in the one thing a transmitter is responsible for.

Two requirements the testbench is strict about:
\begin{itemize}
  \item \textbf{\code{send} is ignored while the transmitter is busy.} Reloading mid-frame would
    restart the byte and corrupt what is already on the line. A request arriving during a frame is
    discarded, not queued.
  \item \textbf{\code{busy} stays high through the whole stop bit}, not just the data bits. The frame
    is not over until the line has been held high for the last bit period.
\end{itemize}

\xpart{d} Run the testbench. It plays receiver: it waits for your start bit, samples the middle of
every bit period, and checks the framing, the byte, the \code{busy} handshake, and that a \code{send}
applied two bit periods into a frame is ignored.

\begin{shell}
cd lectures/L08/exercises/uart_tx8
cp ../../../L04/exercises/reset_sync/reset_sync.vhd .     # the two you wrote yourself
cp ../../../L07/exercises/timer/timer.vhd .
ghdl -a --std=93 reset_sync.vhd timer.vhd uart_tx8.vhd uart_tx8_tb.vhd
ghdl -e --std=93 uart_tx8_tb
ghdl -r --std=93 uart_tx8_tb --assert-level=error --stop-time=10ms
\end{shell}

\xpart{e} Answer, in prose:
\begin{itemize}
  \item A transmitter skipping the stop bit entirely still leaves the line high afterwards, since
    idle and stop are the same level. What actually goes wrong, and when does it first show? The
    testbench catches this one: which of its checks does it, and why could it never have been caught
    by looking at \code{tx} alone?
  \item Your receiver in Exercise~\ref{c8:ex:rx} samples mid-bit; your transmitter changes the line
    at the edges. If both ran on crystals differing by 1~\%, roughly how many bits into a frame would
    the receiver's sampling point drift into the wrong bit? What does that tell you about why serial
    formats have a start bit per byte, rather than a single one at the beginning of a long message?
\end{itemize}

\modulefig[0.66]{lectures/L08/appendix/images/uart_tx8.png}

\begin{selfcheck}
\textbf{Self-check.} Name your entity \code{uart_tx8}, with the generic \code{BIT_TICK_COUNT}
(\code{natural}), the inputs \code{clock}, \code{reset_n}, \code{data_in}
(\code{std_logic_vector(7 downto 0)}), \code{send}, and the outputs \code{tx}, \code{busy}, declared
in that order. Its testbench is in \file{exercises/uart_tx8/}, and nothing else: copy in your own
\code{reset_sync} and \code{timer}, as part \textbf{d)} shows.
\end{selfcheck}
